\section{总体思路}

\subsection{建模主线：把"保证"做成目标函数}

本问的关键在于"第二次测向必须成功、定位必须够准"这两件事都必须是\textbf{保证}而不是
\textbf{期望}。因此本文不采用"误差服从正态分布、最小化均方误差"的统计路线，而采用
\textbf{集员 + 最坏情况}路线：把所有已知信息写成集合（源的不确定集、示向度误差区间），
把"最坏情况下仍然成立"写成约束与目标。这条路线有三个直接好处：结论是硬保证；不需要假设
误差分布；每一步都可以用解析几何独立复核。

建模主线由四步组成（图~\ref{fig:roadmap} 上半支）：

\begin{enumerate}
  \item \textbf{由一次测向写出源不确定集} $C$：以 $S_1$ 为顶点、深 $1500$ m、张角
        $2^\circ$ 的窄楔形（§\ref{sec:model}）；
  \item \textbf{由"一定听得到"写出可行域} $F$：$C$ 中每个源到 $S_2$ 的距离都不超过接收
        半径下界 $1000$ m，等价于 $S_2$ 落在以 $C$ 的四个极点为心、半径 $1000$ m 的四个
        圆盘之交里，是一个面积 $0.447$ km$^2$ 的透镜（§\ref{sec:model}）；
  \item \textbf{由定位区域口径写出代价函数}：对给定的 $S_2$，取源位置与第二次测向误差的
        最坏组合，得到该点的最坏定位直径 $J(S_2)$；
  \item \textbf{求解} $\min_{S_2\in F}J(S_2)$，得到最优第二检测点 $S_2^*$；再以
        $F=J^*/J$ 定义适合度、以 $F\ge1/(1+\eta)$ 定义候选区域，回答"周边区域的适合度"。
\end{enumerate}

\begin{figure}[H]
  \centering
  \begin{tikzpicture}[
      box/.style={draw, rounded corners=2pt, align=center, font=\small, inner sep=3pt,
                  minimum height=1.05cm, text width=3.15cm},
      ar/.style={-{Latex[length=2mm]}, thick},
      node distance=5mm and 6mm
    ]
    \node[box, fill=blue!5]                (c)    {源不确定集 $C$：楔形\\$d\in[5,1500]$ m，$2^\circ$};
    \node[box, fill=blue!5, right=of c]    (f)    {可测性保证：可行域 $F$\\（透镜 $0.447$ km$^2$）};
    \node[box, fill=orange!8, right=of f]  (j)    {最坏定位直径\\$J(S_2)$ 泛函};
    \node[box, fill=orange!14, right=of j] (opt)  {$\min\limits_{S_2\in F}J$：\\$S_2^*=(801,605)$ m};
    \node[box, fill=green!6, below=9mm of c]    (lit)  {文献判据：CRLB /\\GDOP / 几何稀释式};
    \node[box, fill=green!6, below=9mm of f]    (pre)  {GDOP 解析场快筛\\（$5$ m 细网格）};
    \node[box, fill=green!10, below=9mm of j]   (cert) {精确复核 $+$ 细化\\（全局认证）};
    \node[box, fill=gray!10, below=9mm of opt]  (band) {适合度 $F=J^*/J$\\候选弧带 $0.0207$ km$^2$};
    \draw[ar] (c) -- (f);
    \draw[ar] (f) -- (j);
    \draw[ar] (j) -- (opt);
    \draw[ar] (lit) -- (pre);
    \draw[ar] (pre) -- (cert);
    \draw[ar, dashed] (cert) -- (opt);
    \draw[ar, dashed] (opt) -- (band);
    \draw[ar, dashed] (j) -- (band);
  \end{tikzpicture}
  \caption{问题二求解的技术路线。上半支为集员建模与最坏情况最小化的主链；下半支为文献判据层，
  它提供解析快筛并独立认证最优解；右端给出适合度与候选区域。}
  \label{fig:roadmap}
\end{figure}

\subsection{为什么还要引入文献判据}

纯方位测向的交会定位在文献中有成熟的理论刻画：两站测向的 Fisher 信息矩阵、Cram\'er-Rao
下界（CRLB）误差椭圆、几何精度因子（GDOP）以及几何稀释效应。这些工具与本文的集员口径
回答的是同一类问题（"两条带误差的射线交会有多准"），但语汇不同：文献口径是\textbf{统计}
尺度（独立误差平方向加），本文口径是\textbf{有界}尺度（误差可同向叠加）。二者是否一致、
能否互相替代，是审稿与工程上都必须回答的问题。

本文的处理方式是\textbf{把文献判据实现成一层代码，逐项与精确判据对照}（§\ref{sec:lit}）：

\begin{itemize}
  \item 文献的 CRLB 椭圆主半轴、GDOP、几何稀释式都能写成 $1/\sin\gamma$ 主导的闭式
        （$\gamma$ 为源处两条示向度的交会角），这与本文一阶解析式同源，说明\textbf{趋势一致}；
  \item 但在本模型最坏情形上，文献式比精确最坏半径\textbf{低约三成}（$46.98$ m 对
        $67.04$ m），因为源不确定集取最坏、误差同向叠加这两条本文必须承担，
        而统计口径把它们平均掉了。清除判据是 $20$ m 硬阈值，低估三成的量不能当保险界；
  \item 文献口径与精确判据的\textbf{排序几乎一致}（可行域内 $1276$ 个采样点上
        Spearman 秩相关 $0.998$），因此可以拿它做\textbf{快速预筛}：用解析 GDOP 场在 $5$ m
        细网格上筛出少数候选，再用精确构造复核并细化。
\end{itemize}

这一步的结果不是"换一个更好的判据"，而是\textbf{换来了一个独立的全局认证}：细网格
快筛路线给出的最优解与 $25$ m 粗搜 $+$ 局部细化路线相差仅 $8.5$ m、最坏直径优
$0.063$ m，两条路线互相印证，排除了"粗网格漏掉更好盆地"的风险。此外，文献常用的
\textbf{期望}精度口径（最小化平均定位直径）与本文的 minimax 口径给出的最优点相距不足
$20$ m、最坏直径相差小于 $0.1\%$，说明最终推荐区域\textbf{对判据的选择不敏感}——
这是本文结论稳健性的直接证据（§\ref{sec:res}）。

\subsection{全文章节安排}

§\ref{sec:model} 建立模型（不确定集、可行域、代价泛函、适合度与候选区域、一阶解析式）；
§\ref{sec:lit} 建立文献判据层并逐项对照；§\ref{sec:algo} 给出算法（离散化、粗搜、细化、
文献快筛与全局认证、候选弧带、镜像规范）；§\ref{sec:res} 给出求解结果与两张成果图；
§\ref{sec:sens} 给出灵敏度分析；§\ref{sec:eval} 给出模型校验、评价与推广。
